\begin{definition} 
    $\PATH = \{(G, s, t) \ | \text{ в орграфе $G$ есть путь из $s$ в $t$}\}$.
\end{definition}

\begin{theorem} 
    $\PATH \in \nlcsp$.
\end{theorem}

\begin{proof}
    \begin{enumerate}
        \item $\PATH \in \nlsp$. Воспользуемся сертификатным определением, где сертификатом будет путь. Теперь верификатор помнит последнюю вершину, читает следующую и проверяет, есть ли ребро. В итоге память логарифмическая.
        \item Докажем полноту. Построим конфигурационный граф машины. Принятие $x$ для некоторого языка и его МТ означает, что есть путь из стартового состояния в завершающее. Собственно, $\PATH$ проверяет, есть ли такой путь. То есть стартовое состояние соответствует $s$, а принимающее — $t$. Теперь опишем логарифмическую сводимость. То есть надо уметь на логарифмической памяти понимать, можно ли перейти из данного состояния в следующее и соответствует ли это программе. Первое равносильно тому, что лента изменилась лишь в том месте, куда указывал указатель, а второе заключается в сравнении нового состояния со всеми доступными по программе, для чего нам хватит логарифмической памяти.
    \end{enumerate}
\end{proof}

\begin{definition}
    $\mathsf{SCONNECTED}=\{G|$ из любой вершины ориентированного графа $G$ в любую другую идет ориентированный путь\} (сильно связные графы)
\end{definition}

\begin{theorem}
    Язык $\mathsf{SCONNECTED}$ является $\nlsp$-полным
\end{theorem} 

\begin{proof}
    покажем, что $\mathsf{SCONNECTED}\in\nlsp$. В качестве сертификата будем передавать
    \[((1, 1),1\xrightarrow{path_{1,1}}1), ((1, 2),1\xrightarrow{path_{1,2}}2),\ldots, ((1, n),1\xrightarrow{path_{1,n}}n), ((2, 1),2\xrightarrow{path_{2,1}}1),\ldots,((n, n),n\xrightarrow{path_{n,n}}n)\]
    
    Проверка будет такая: перебираем $i=1\ldots n$, внутри перебираем $j=1\ldots n$ 
    внутри проверяем, что есть путь от $i$ до $j$, используя сертификат.
    
    Сведём $\mathsf{PATH}$ к $\mathsf{SCONNECTED}$. Пусть есть $(G,s,t)$ преобразуем $G$ добавив для любой вершины $v$ ориентированное ребро из $v$ в $s$ (кроме $v=s$). Далее для любой вершины $v$ ориентированное ребро из $t$ в $v$ (кроме $v=t$). Полученный граф обозначим $G'$. Заметим, что функция $G\to G'$ логарифмически вычислимая.
    
    Если есть путь из $s$ в $t$ ($s\xrightarrow{p}t$) в графе $G$, то $\forall u,v$ есть путь $u\rightarrow s\xrightarrow{p}t\rightarrow v$. Значит $G'$ будет лежать в $\mathsf{SCONNECTED}$.
    
    Обратно, пусть $G'$ сильно связный, тогда есть путь из $s$ в $t$ в $G'$, но этот путь можно преобразовать так, чтобы он не использовал добавленные ребра, значит есть путь в $G$ из $s$ в $t$, то есть $(G,s,t)\in\mathsf{PATH}$.
    
    Получили
    \[(G,s,t)\in\mathsf{PATH}\iff G' \in \mathsf{SCONNECTED}\]
    
    Значит $\mathsf{SCONNECTED}$ $\nlsp$-полный.
\end{proof} 